% -------------------------------------------------------------------------
% WBS ID: RES-NUM-SPEC
% Deliverable: Integrated Numerical-Method Specification
% Status: Draft complete
% Evidence status: Regional and local URANS--VOF research specifications complete
% Review status: Not reviewed
% Last updated: 2026-07-18
% Dependencies: RES-NUM-FRM, RES-NUM-MSH, RES-NUM-IBC, RES-NUM-TIM, RES-NUM-STB, RES-NUM-ERR
% Downstream interfaces: RES-VER-FRM, RES-VER-BMK, RES-VER-MET, Software local CFD implementation
% -------------------------------------------------------------------------

\section{RES-NUM-SPEC --- Integrated Numerical-Method Specification}
\label{sec:res-num-spec}

\subsection{Purpose and Scope}
\label{sec:res-num-spec-purpose}

This Deliverable records the integrated numerical decision status across
the regional hydrodynamic solver, local three-dimensional CFD solver,
structural solver, two-dimensional--three-dimensional coupling solver,
and FSI coupling solver.

It does not mark the whole numerical framework as complete. The
regional nonlinear shallow-water finite-volume branch is now complete
at research-specification level. The local three-dimensional
URANS--VOF hydrodynamic branch is now specified at solving-procedure
level. Structural dynamics, regional-to-local transfer details, FSI
coupling, implementation and validation remain separate workstreams.

\subsection{Research Questions}
\label{sec:res-num-spec-research-questions}

The current phase records:

\begin{enumerate}
  \item which solver-family decisions are complete;
  \item which lower-level regional and local choices remain provisional
    or unresolved;
  \item which structural and coupling solver choices remain open;
  \item which software architecture references are selected without
    constraining every solver to one implementation.
\end{enumerate}

\subsection{Terminology and Definitions}
\label{sec:res-num-spec-definitions}

The decision statuses are:

\begin{description}
  \item[selected] The project adopts the item for the current branch.
  \item[provisional] The item is the current working direction but still
    requires lower-level comparison or verification.
  \item[excluded] The item is removed from the active branch.
  \item[unresolved] The item remains open and requires later evidence.
\end{description}

\subsection{Breadth-First Numerical Research Strategy}
\label{sec:res-num-spec-breadth-first-strategy}

The project adopts a breadth-first numerical research strategy:

\begin{quote}
  define the governing models
  $\rightarrow$
  select a numerical solver family for each model
  $\rightarrow$
  select numerical coupling methods
  $\rightarrow$
  return to shared framework components
  $\rightarrow$
  increase detail within each selected branch
\end{quote}

The current solver-family order is:

\begin{enumerate}
  \item regional NLSWE solver;
  \item local three-dimensional two-phase fluid solver;
  \item structural dynamics solver;
  \item two-dimensional--three-dimensional coupling solver;
  \item fluid--structure coupling solver.
\end{enumerate}

After these family-level decisions are established, the research shall
return to shared numerical components: initial conditions, boundary
conditions, source-term discretisation, mesh and refinement, numerical
fluxes, reconstruction, time integration, linear and nonlinear solvers,
coupling residuals, convergence criteria, verification, and validation.

The finite-volume selection in `RES-NUM-FRM` completes the regional
solver-family decision. The local three-dimensional selection in
\cref{sec:num-local-three-dimensional-fluid-solver-selection} completes
the local fluid-solver family decision while leaving lower-level local
algorithms provisional or unresolved.

\subsection{Numerical Decision Status}
\label{sec:res-num-spec-numerical-decision-status}

\begin{table}[htbp]
  \centering
  \small
  \renewcommand{\arraystretch}{1.2}
  \begin{tabular}{
      p{0.18\textwidth}
      p{0.14\textwidth}
      p{0.48\textwidth}
    }
    \hline
    Item
    &
    Status
    &
    Decision statement
    \\
    \hline
    Regional governing model
    &
    selected
    &
    Two-dimensional nonlinear shallow-water equations with variable
    bathymetry, source forcing, wetting and drying, and optional later
    dispersive extension
    \\
    Regional numerical family
    &
    selected
    &
    Cell-centred finite-volume method, as recorded in
    \cref{sec:res-num-frm-regional-solver-decision}
    \\
    Regional solver composition
    &
    selected
    &
    Polygonal finite-volume update with triangular and quadrilateral
    practical meshes, MUSCL reconstruction, weighted least-squares
    gradients, Barth--Jespersen limiter, hydrostatic reconstruction,
    HLLC primary flux with HLL/HLLE fallback, wet/dry positivity
    controls, SSPRK(3,3) transport, Strang-split Manning friction,
    CFL timestep, open/radiation boundaries and optional sponge damping
    \\
    Finite difference for production regional solver
    &
    excluded
    &
    Not selected for production, retained as a numerical reference and
    source of semi-implicit architectural ideas
    \\
    Lattice/discrete Boltzmann for regional solver
    &
    excluded
    &
    Removed from the active shortlist pending tsunami validation,
    dry-front, positivity, well-balancing and complex-bathymetry evidence
    \\
    Spectral and pseudospectral methods
    &
    excluded
    &
    Screened out for this project branch owing to irregular coastal
    boundaries, wetting and drying, bores, discontinuities, bathymetry
    and local refinement needs
    \\
    OpenFOAM local CFD reference
    &
    selected
    &
    Principal open-source software and architectural reference for the
    local three-dimensional CFD model
    \\
    Local three-dimensional numerical family
    &
    selected
    &
    Eulerian cell-centred polyhedral finite-volume method with
    volume-of-fluid interface capture, pressure-correction or projection
    incompressibility treatment, URANS $k$--$\omega$ SST operational
    closure, and LES retained as a higher-fidelity comparison mode
    \\
    Local three-dimensional solve procedure
    &
    selected
    &
    Segregated URANS--VOF update with regional inlet forcing, timestep
    selection, bounded VOF transport, mixture-property update,
    momentum predictor, pressure correction, turbulence solve,
    wall-function update, metric extraction and accept/reject control
    \\
    Structural solver
    &
    unresolved
    &
    Element family, constitutive discretisation and structural time
    integration remain to be selected
    \\
    Two-dimensional--three-dimensional numerical coupling
    &
    unresolved
    &
    Transfer operators, coupling residuals and convergence criteria
    remain to be selected
    \\
    FSI numerical coupling
    &
    unresolved
    &
    One-way load transfer is retained and two-way coupling remains
    conditional on structural feedback; the numerical algorithm is not
    yet selected
    \\
    \hline
  \end{tabular}
  \caption{Current numerical decision status for the breadth-first
  research phase}
  \label{tab:res-num-spec-numerical-decision-status}
\end{table}

\subsection{Literature Evidence}
\label{sec:res-num-spec-literature-evidence}

% TODO: Record evidence from peer-reviewed papers, authoritative datasets,
% standards, technical reports and primary documentation.
%
% Do not create a detached literature-summary list. Explain how each source
% supports, contradicts or limits a project decision.

\subsection{Governing Relations and Quantitative Definitions}
\label{sec:res-num-spec-governing-relations}

% TODO: Record relevant equations, dimensionless groups, empirical models,
% extraction rules or quantitative definitions.
%
% Use align environments for displayed equations.
% Every displayed equation must have a unique \label{eq:...}.
% Do not place punctuation after displayed equations.

\subsection{Discrete Formulation}
\label{sec:res-num-spec-discrete-formulation}

The integrated numerical framework shall permit solver families to
differ across governing systems. The regional solver may use a
standalone finite-volume implementation, while the local
three-dimensional CFD solver may use OpenFOAM as software and
architecture reference.

The regional implementation shall not be forced inside OpenFOAM. Shared
abstractions shall be transferred where they improve consistency across
solvers.

\subsubsection{Regional 2D NLSWE Finite-Volume Specification}
\label{sec:res-num-spec-regional-2d-nlswe-specification}

The regional 2D NLSWE numerical model is established at
research-specification level. The solved state is the horizontal
two-dimensional conserved vector:

\begin{align}
  \mathbf{U}_i(t)
  &=
  \frac{1}{A_i}
  \int_{K_i}
  \mathbf{U}(\mathbf{x},t)
  \,\dd A
  \label{eq:res-num-spec-regional-cell-average}
\end{align}

For cell $K_i$ with area $A_i$, boundary faces $f$, face lengths
$L_f$, outward normals $\mathbf{n}_f$, and source contribution
$\mathbf{S}_i$, the finite-volume update is:

\begin{align}
  \frac{\dd\mathbf{U}_i}{\dd t}
  &=
  -
  \frac{1}{A_i}
  \sum_{f\in\partial K_i}
  L_f
  \widehat{\mathbf{F}}_{if}
  +
  \mathbf{S}_i
  \label{eq:res-num-spec-regional-fv-update}
\end{align}

For a polygon with $N_i$ faces this is written:

\begin{align}
  \frac{\dd\mathbf{U}_i}{\dd t}
  &=
  -
  \frac{1}{A_i}
  \sum_{k=1}^{N_i}
  L_{i,k}
  \widehat{\mathbf{F}}_{i,k}
  +
  \mathbf{S}_i
  \label{eq:res-num-spec-regional-polygon-update}
\end{align}

Cell shape changes only the geometry entering the update:

\begin{align}
  A_i,\quad
  N_i,\quad
  L_{i,k},\quad
  \mathbf{n}_{i,k},\quad
  \mathbf{x}_{f_{i,k}},\quad
  \mathcal{N}(i)
  \label{eq:res-num-spec-regional-cell-geometry-set}
\end{align}

Primary practical regional meshes use unstructured triangles. Quads are
retained for tests and simple offshore corridors. General straight-edged
polygons remain admissible where mesh quality and validation evidence
support them.

\paragraph{Geometry.}
\label{sec:res-num-spec-regional-geometry}

For counter-clockwise vertices
$\mathbf{x}_{i,1},\ldots,\mathbf{x}_{i,N_i}$, with
$\mathbf{x}_{i,N_i+1}=\mathbf{x}_{i,1}$:

\begin{align}
  L_{i,k}
  &=
  \left\|
    \mathbf{x}_{i,k+1}
    -
    \mathbf{x}_{i,k}
  \right\|
  \label{eq:res-num-spec-regional-face-length}
  \\
  \mathbf{x}_{f_{i,k}}
  &=
  \frac{
    \mathbf{x}_{i,k}
    +
    \mathbf{x}_{i,k+1}
  }{2}
  \label{eq:res-num-spec-regional-face-centroid}
  \\
  \mathbf{t}_{i,k}
  &=
  \frac{
    \mathbf{x}_{i,k+1}
    -
    \mathbf{x}_{i,k}
  }{L_{i,k}}
  \label{eq:res-num-spec-regional-face-tangent}
  \\
  \mathbf{n}_{i,k}
  &=
  \left[
    t_y,\,
    -t_x
  \right]^{\mathsf{T}}
  \label{eq:res-num-spec-regional-face-normal}
\end{align}

For a triangular cell:

\begin{align}
  A_i
  &=
  \frac{1}{2}
  \left|
    \left(
      \mathbf{x}_2-\mathbf{x}_1
    \right)
    \times
    \left(
      \mathbf{x}_3-\mathbf{x}_1
    \right)
  \right|
  \label{eq:res-num-spec-regional-triangle-area}
  \\
  \mathbf{x}_i
  &=
  \frac{
    \mathbf{x}_1
    +
    \mathbf{x}_2
    +
    \mathbf{x}_3
  }{3}
  \label{eq:res-num-spec-regional-triangle-centroid}
\end{align}

For a counter-clockwise quadrilateral:

\begin{align}
  J_k
  &=
  x_k y_{k+1}
  -
  x_{k+1} y_k
  \label{eq:res-num-spec-regional-quadrilateral-jacobian-term}
  \\
  A_i
  &=
  \frac{1}{2}
  \sum_{k=1}^{4}
  J_k
  \label{eq:res-num-spec-regional-quadrilateral-area}
  \\
  x_i
  &=
  \frac{1}{6A_i}
  \sum_{k=1}^{4}
  \left(
    x_k+x_{k+1}
  \right)
  J_k
  \label{eq:res-num-spec-regional-quadrilateral-centroid-x}
  \\
  y_i
  &=
  \frac{1}{6A_i}
  \sum_{k=1}^{4}
  \left(
    y_k+y_{k+1}
  \right)
  J_k
  \label{eq:res-num-spec-regional-quadrilateral-centroid-y}
\end{align}

\paragraph{Reconstruction and Flux.}
\label{sec:res-num-spec-regional-reconstruction-flux}

The reconstructed variables are:

\begin{align}
  \mathbf{W}
  &=
  \left[
    \eta,\,
    u,\,
    v,\,
    b
  \right]^{\mathsf{T}}
  \label{eq:res-num-spec-regional-reconstruction-state}
\end{align}

For scalar $\phi$, the face value is reconstructed as:

\begin{align}
  \phi_{i,f}
  &=
  \phi_i
  +
  \alpha_i^{\phi}
  \nabla\phi_i
  \cdot
  \left(
    \mathbf{x}_f-\mathbf{x}_i
  \right)
  \label{eq:res-num-spec-regional-muscl-reconstruction}
\end{align}

Weighted least-squares gradients are selected:

\begin{align}
  \mathbf{d}_{ij}
  &=
  \mathbf{x}_j
  -
  \mathbf{x}_i
  \label{eq:res-num-spec-regional-neighbour-displacement}
  \\
  \mathbf{M}_i
  &=
  \sum_{j\in\mathcal{N}(i)}
  w_{ij}
  \mathbf{d}_{ij}
  \mathbf{d}_{ij}^{\mathsf{T}}
  \label{eq:res-num-spec-regional-ls-matrix}
  \\
  \mathbf{r}_i
  &=
  \sum_{j\in\mathcal{N}(i)}
  w_{ij}
  \mathbf{d}_{ij}
  \left(
    \phi_j-\phi_i
  \right)
  \label{eq:res-num-spec-regional-ls-rhs}
  \\
  \nabla\phi_i
  &=
  \mathbf{M}_i^{-1}
  \mathbf{r}_i
  \label{eq:res-num-spec-regional-ls-gradient}
\end{align}

The baseline limiter is Barth--Jespersen. Reconstruction is reverted to
first order near wet/dry fronts where required for positivity.

HLLC is the primary face flux for wet, valid states. HLL or HLLE is the
fallback near dry, degenerate or unsafe states. Each face flux is solved
in normal/tangential coordinates:

\begin{align}
  u_n
  &=
  u n_x
  +
  v n_y
  \label{eq:res-num-spec-regional-normal-velocity}
  \\
  u_t
  &=
  -u n_y
  +
  v n_x
  \label{eq:res-num-spec-regional-tangential-velocity}
  \\
  \widetilde{\mathbf{U}}_K
  &=
  \left[
    h_K,\,
    h_K u_{n,K},\,
    h_K u_{t,K}
  \right]^{\mathsf{T}},
  \qquad
  K\in\{L,R\}
  \label{eq:res-num-spec-regional-rotated-state}
\end{align}

\paragraph{Hydrostatic Reconstruction and Sources.}
\label{sec:res-num-spec-regional-hydrostatic-source}

Hydrostatic reconstruction uses:

\begin{align}
  b_f^{\mathrm{H}}
  &=
  \max
  \left(
    b_L,\,
    b_R
  \right)
  \label{eq:res-num-spec-regional-hydrostatic-face-bed}
  \\
  h_L^{\mathrm{H}}
  &=
  \max
  \left(
    0,\,
    \eta_L-b_f^{\mathrm{H}}
  \right)
  \label{eq:res-num-spec-regional-hydrostatic-left-depth}
  \\
  h_R^{\mathrm{H}}
  &=
  \max
  \left(
    0,\,
    \eta_R-b_f^{\mathrm{H}}
  \right)
  \label{eq:res-num-spec-regional-hydrostatic-right-depth}
\end{align}

The second-order bed source is:

\begin{align}
  \mathbf{S}_{\mathrm{b},i}^{(2)}
  &=
  \left[
    0,\,
    -g h_i (b_x)_i,\,
    -g h_i (b_y)_i
  \right]^{\mathsf{T}}
  \label{eq:res-num-spec-regional-second-order-bed-source}
\end{align}

The sponge source uses:

\begin{align}
  \mathbf{S}_{\mathrm{sp},i}
  &=
  -\sigma_i
  \left[
    h_i-h_{\mathrm{ref},i},\,
    q_{x,i},\,
    q_{y,i}
  \right]^{\mathsf{T}}
  \label{eq:res-num-spec-regional-sponge-source}
  \\
  h_{\mathrm{ref},i}
  &=
  \max
  \left(
    \eta_{\mathrm{ref}}
    -
    b_i,\,
    0
  \right)
  \label{eq:res-num-spec-regional-sponge-reference-depth}
  \\
  \sigma_i
  &=
  \sigma_{\max}
  \left(
    \frac{d_i}{L_{\mathrm{sp}}}
  \right)^p
  \label{eq:res-num-spec-regional-sponge-ramp}
  \\
  \sigma_{\max}
  &=
  \frac{(p+1)c_{\mathrm{ref}}}{L_{\mathrm{sp}}}
  \ln
  \left(
    \frac{1}{A_{\mathrm{target}}}
  \right)
  \label{eq:res-num-spec-regional-sponge-max}
\end{align}

The baseline sponge values are $p=2$, $A_{\mathrm{target}}=0.01$, and
$L_{\mathrm{sp}}=50$--$100~\mathrm{km}$ before adjustment to the final
domain scale.

\paragraph{Wetting, Drying and Time Integration.}
\label{sec:res-num-spec-regional-wet-dry-time}

The regional solver uses positivity-preserving wet/dry treatment, not a
regional VOF method. Dry cells are reset by:

\begin{align}
  h_i
  <
  h_{\mathrm{dry}}
  \quad
  &\Rightarrow
  \quad
  h_i=0,\quad
  q_{x,i}=q_{y,i}=0
  \label{eq:res-num-spec-regional-dry-reset}
\end{align}

The draining limiter is:

\begin{align}
  \lambda_i
  &=
  \min
  \left(
    1,\,
    \frac{
      A_i
      \max
      \left(
        h_i-h_{\mathrm{dry}},\,
        0
      \right)
    }{
      \Delta t
      \sum_f
      L_f
      \max
      \left(
        F_{h,if},\,
        0
      \right)
      +
      \varepsilon
    }
  \right)
  \label{eq:res-num-spec-regional-draining-limiter}
\end{align}

The transport/friction time update is:

\begin{align}
  \mathbf{U}^{n+1}
  &=
  \mathcal{M}_{\Delta t/2}
  \circ
  \mathcal{T}_{\Delta t}^{\mathrm{SSPRK3}}
  \circ
  \mathcal{M}_{\Delta t/2}
  \left(
    \mathbf{U}^{n}
  \right)
  \label{eq:res-num-spec-regional-strang-ssprk3}
\end{align}

The CFL timestep is:

\begin{align}
  a_f
  &=
  \max_K
  \left(
    \left|
      u_{n,K}
    \right|
    +
    \sqrt{g h_K}
  \right)
  \label{eq:res-num-spec-regional-face-signal-speed}
  \\
  \Delta t_i
  &=
  C_{\mathrm{CFL}}
  \frac{A_i}{
    \sum_{f\in\partial K_i}
    L_f a_f
  }
  \label{eq:res-num-spec-regional-cell-timestep}
  \\
  \Delta t
  &=
  \min_i
  \Delta t_i
  \label{eq:res-num-spec-regional-global-timestep}
\end{align}

\paragraph{Established and Deferred Scope.}
\label{sec:res-num-spec-regional-established-deferred}

\begin{table}[htbp]
  \centering
  \small
  \renewcommand{\arraystretch}{1.2}
  \begin{tabular}{
      p{0.42\textwidth}
      p{0.42\textwidth}
    }
    \hline
    Established for the regional model
    &
    Deferred or separate
    \\
    \hline
    Regional 2D NLSWE finite-volume formulation; polygonal geometry;
    triangle/quadrilateral practical meshes; reconstruction and limiter;
    hydrostatic reconstruction; HLLC/HLL flux strategy; wetting/drying;
    dynamic bed treatment; SSPRK(3,3) transport with split Manning;
    CFL timestep; open/radiation and sponge treatment; mesh-refinement
    principles
    &
    QGIS/geospatial domain lock; exact bathymetry/topography dataset;
    exact fault-source dataset; validation runs; local 3D/interface/FSI
    model; production software implementation; optimisation and ML
    workflow
    \\
    \hline
  \end{tabular}
  \caption{Regional numerical-model established/deferred status}
  \label{tab:res-num-spec-regional-established-deferred}
\end{table}

\subsection{Conservation Properties}
\label{sec:res-num-spec-conservation-properties}

% TODO: Complete this RES-NUM Domain-specific subsection.

\subsection{Consistency and Formal Accuracy}
\label{sec:res-num-spec-consistency-formal-accuracy}

% TODO: Complete this RES-NUM Domain-specific subsection.

\subsection{Stability Requirements}
\label{sec:res-num-spec-stability-requirements}

% TODO: Complete this RES-NUM Domain-specific subsection.

\subsection{Mesh and Time-Step Requirements}
\label{sec:res-num-spec-mesh-time-step-requirements}

% TODO: Complete this RES-NUM Domain-specific subsection.

\subsection{Numerical Failure Modes}
\label{sec:res-num-spec-numerical-failure-modes}

% TODO: Complete this RES-NUM Domain-specific subsection.

\subsection{Verification Requirements}
\label{sec:res-num-spec-verification-requirements}

% TODO: Complete this RES-NUM Domain-specific subsection.

\subsection{Candidate Approaches}
\label{sec:res-num-spec-candidate-approaches}

The active breadth-first branches are:

\begin{itemize}
  \item regional NLSWE finite-volume solver;
  \item local three-dimensional two-phase CFD solver;
  \item structural dynamics solver;
  \item two-dimensional--three-dimensional coupling solver;
  \item fluid--structure coupling solver.
\end{itemize}

\subsection{Critical Comparison}
\label{sec:res-num-spec-critical-comparison}

% TODO: Compare candidates using physically and project-relevant criteria.
% Do not select an approach solely on popularity or implementation ease.

\subsection{Assumptions and Applicability}
\label{sec:res-num-spec-assumptions}

% TODO: State assumptions and the conditions under which they are valid.

\subsection{Limitations and Failure Conditions}
\label{sec:res-num-spec-limitations}

% TODO: State where the evidence, model or selected approach ceases to be
% reliable or applicable.

\subsection{Project-Specific Conclusions}
\label{sec:res-num-spec-conclusions}

% TODO: Convert the evidence into conclusions specific to the Studentship
% project. Do not merely restate literature findings.

The project has selected the regional finite-volume family, completed
the regional 2D NLSWE numerical model at research-specification level,
and selected the local three-dimensional two-phase finite-volume solver
family. It has not completed structural, dimensional-coupling or FSI
numerical-method selection. The numerical architecture shall remain
modular while implementation and validation work proceeds.

\subsection{Selected Approach and Rationale}
\label{sec:res-num-spec-selected-approach}

% TODO: Record the selected approach only after sufficient evidence exists.
% State the alternatives rejected and the reasons for rejection.

\subsubsection{Selection of the Local Three-Dimensional Fluid Solver}
\label{sec:num-local-three-dimensional-fluid-solver-selection}

The selected numerical framework for the local three-dimensional fluid
model is an Eulerian, cell-centred, polyhedral finite-volume method
combined with a volume-of-fluid interface-capturing method and a
pressure-correction procedure. This framework will discretise the
incompressible, immiscible, two-phase Navier--Stokes equations. The same
fluid-solver architecture will support either unsteady
Reynolds-averaged Navier--Stokes or large-eddy simulation turbulence
treatment, with URANS $k$--$\omega$ SST selected as the operational
baseline and LES retained for selected high-fidelity comparison cases.

The numerical-family decision is therefore:

\begin{quote}
  \centering
  cell-centred polyhedral finite volume
  $+$ volume-of-fluid interface capture
  $+$ pressure correction
\end{quote}

This decision fixes the principal spatial, interface and
incompressibility enforcement families. It does not yet fix the exact
volume-of-fluid algorithm, pressure-correction algorithm, temporal
discretisation, wall-function constants, mesh-motion method or
fluid--structure coupling procedure.

\paragraph{Conservative Fluid and Phase Transport.}
\label{sec:num-local-three-dimensional-conservation-rationale}

The finite-volume method evaluates the change in each control-volume
state through fluxes crossing its bounding faces. This provides a direct
conservative representation of fluid transport and supplies a consistent
basis for calculating momentum exchange and structural loading. A common
face flux also enables neighbouring control volumes to exchange equal
and opposite contributions, which is advantageous for global
conservation auditing and domain decomposition.

The volume-of-fluid method was selected as the free-surface
representation since it tracks the volume fraction of each phase and is
therefore naturally compatible with conservative finite-volume
transport. Geometric volume-of-fluid methods have been developed for
arbitrary polyhedral meshes and have demonstrated strong phase-volume
conservation, boundedness and interface sharpness
\parencite{roenby2016computational}.

\paragraph{Breaking, Overtopping and Interface Topology.}
\label{sec:num-local-three-dimensional-interface-topology-rationale}

The local solver must represent wave breaking, overturning, overtopping,
jet formation, fragmentation and interface reconnection. An Eulerian
interface-capturing method allows these changes in topology without
requiring the computational mesh to conform to the air--water interface.

A sharp finite-volume free-surface formulation combining geometric
volume-of-fluid advection with a Ghost Fluid treatment has been verified
for gravity-wave propagation and validated for violent green-water flow
\parencite{vukcevic2018sharp}. This provides direct methodological
evidence that the selected family can represent both smooth wave
propagation and strongly deformed free-surface flow within one solver
architecture.

\paragraph{Pressure and Structural-Load Prediction.}
\label{sec:num-local-three-dimensional-pressure-load-rationale}

The structural analysis requires reliable predictions of local pressure,
pressure impulse, resultant force, shear traction and overturning
moment. The pressure-correction method enforces the incompressibility
constraint by coupling the momentum predictor to a pressure equation and
subsequently correcting the velocity and face-flux fields.

The large density change across the air--water interface requires
particular care. A diffuse pressure and density transition may alter the
magnitude and timing of an impact load. Sharp-interface finite-volume
formulations provide a route for imposing the pressure-gradient and
density discontinuities consistently while retaining a compact stencil
on general cell shapes \parencite{vukcevic2018sharp}.

\paragraph{Mesh and Geometric Flexibility.}
\label{sec:num-local-three-dimensional-mesh-rationale}

Polyhedral finite-volume meshes permit body-fitted representation of
irregular barriers, coastal topography and local three-dimensional
structures. They also allow the mesh resolution to be concentrated
around the free surface, barrier edges, overtopping regions and wakes
without imposing a uniformly fine Cartesian grid.

Geometric volume-of-fluid methods have been formulated for structured
and unstructured polyhedral meshes \parencite{roenby2016computational}.
A related OpenFOAM implementation has demonstrated unstructured
geometric volume of fluid with local dynamic adaptive mesh refinement
\parencite{maric2013vofoam}. The exact refinement algorithm remains
unresolved and will be selected during the detailed mesh-framework
study.

\paragraph{URANS and LES Compatibility.}
\label{sec:num-local-three-dimensional-turbulence-compatibility}

Unsteady Reynolds-averaged Navier--Stokes and large-eddy simulation will
be treated as alternative operating modes of the same local
three-dimensional solver. Both modes will use the same continuity,
momentum, phase-fraction, pressure-correction, boundary and
force-integration infrastructure.

The unsteady Reynolds-averaged Navier--Stokes mode will provide the
operational framework for repeated barrier assessment and optimisation.
The large-eddy simulation mode will provide higher-fidelity analysis of
selected impact cases in which resolved separation, breaking-generated
turbulence, fluctuating pressure or load spectra are important. The
baseline Reynolds-averaged closure is $k$--$\omega$ SST
\parencite{menter1994twoequation}. The LES subgrid-scale closure,
rough-wall treatment and wall-function constants remain unresolved.

\paragraph{Local URANS--VOF Solving Procedure.}
\label{sec:num-local-three-dimensional-solving-procedure}

At accepted time level $t^n$, the known local fields are:

\begin{align}
  \mathcal{Q}^{n}_{\mathrm{local}}
  &=
  \left\{
    \alpha^n,\,
    \rho^n,\,
    \mu^n,\,
    \mathbf{u}^n,\,
    p^n,\,
    k^n,\,
    \omega^n,\,
    \mu_t^n,\,
    \phi^n
  \right\}
  \label{eq:res-num-spec-local-known-fields}
\end{align}

The local segregated update is:

\begin{enumerate}
  \setcounter{enumi}{-1}
  \item Interpolate the accepted regional two-dimensional forcing to
    the local three-dimensional inlet using the selected
    regional-to-local history and the boundary reconstruction in
    `RES-MOD-IBC`.
  \item Select the timestep from the flow-CFL, interface-CFL,
    gravity-wave and diffusion constraints in
    \cref{eq:res-num-tim-local-timestep-minimum}.
  \item Update all boundary-face states for
    $\alpha$, $\rho$, $\mu$, $\mathbf{u}$, $p$, $k$, $\omega$ and
    $\phi$.
  \item Solve the bounded compressive VOF equation for
    $\alpha^{n+1}$.
  \item Update $\rho^{n+1}$ and $\mu^{n+1}$ from
    $\alpha^{n+1}$ using the mixture-property definitions in
    `RES-MOD-EQS`.
  \item Assemble and solve the URANS momentum predictor for
    $\mathbf{u}^{*}$.
  \item Construct pressure-consistent face fluxes $\phi^{*}$ using
    Rhie--Chow or an equivalent collocated-mesh interpolation.
  \item Solve the pressure-correction equation for $p'$.
  \item Correct pressure, velocity and face fluxes.
  \item Repeat pressure correctors until the mass residual in
    \cref{eq:res-num-stb-local-mass-residual} is acceptable.
  \item Solve the $k$ and $\omega$ equations using the corrected
    velocity and corrected face flux.
  \item Update $\mu_t^{n+1}$ and $\mu_{\mathrm{eff}}^{n+1}$.
  \item Apply scalable wall-function quantities on terrain and barrier
    wall patches.
  \item Extract pressure, shear, force, moment, run-up, overtopping,
    transmission and energy metrics.
  \item Check boundedness, residual, conservation, mesh-quality and
    output-quality acceptance criteria.
  \item Accept the fields and advance to the next timestep, or reject
    the timestep and repeat with a reduced $\Delta t$.
\end{enumerate}

For any local transported equation with a linearised source
$S_{\psi,i}=S_{C,i}+S_{P,i}\psi_i$, the coefficient updates are:

\begin{align}
  a_{P,i}^{\mathrm{new}}
  &=
  a_{P,i}^{\mathrm{base}}
  -
  V_i S_{P,i}
  \label{eq:res-num-spec-local-source-ap-update}
  \\
  b_i^{\mathrm{new}}
  &=
  b_i^{\mathrm{base}}
  +
  V_i S_{C,i}
  \label{eq:res-num-spec-local-source-b-update}
\end{align}

Pressure correction follows the operator-splitting principle of
\textcite{issa1986operator}. VOF transport follows the source role
`3D-VOF`, with \textcite{hirt1981vof},
\textcite{roenby2016computational} and \textcite{vukcevic2018sharp}
supporting the phase-fraction and sharp-interface requirements.
Wave-boundary state updates use source role `3D-WAVEBC`, supported by
\textcite{higuera2013realistic} and
\textcite{jacobsen2012waves2foam}. Project-specific conclusion: the
solving procedure is event-independent; event-specific inputs change
the inlet history, terrain, bathymetry and geometry data, not the
algorithmic ordering.

\paragraph{Compatibility with the Project Software Architecture.}
\label{sec:num-local-three-dimensional-software-compatibility}

The selected framework is directly aligned with the decision to use
OpenFOAM as the principal open-source software and architectural
reference for the local three-dimensional model. The relevant shared
abstractions include control volumes, owner--neighbour face
connectivity, face-normal fluxes, boundary patches, field operations,
pressure--velocity coupling, two-phase interface transport,
turbulence-model modularity, domain decomposition, solver dictionaries,
runtime configuration and source-level extensibility.

This alignment does not require the regional and local solvers to share
one executable. It provides a consistent numerical vocabulary across the
regional finite-volume solver, the local OpenFOAM-informed solver and
the two-dimensional--three-dimensional coupling interface. The
finite-volume method remains selected on physical and numerical grounds;
software compatibility provides an additional implementation advantage
rather than the primary justification.

\paragraph{Assessment of Alternative Numerical Families.}
\label{sec:num-local-three-dimensional-alternative-assessment}

Stabilised finite elements remain a credible alternative for complex
geometry, variational multiscale turbulence modelling and monolithic
fluid--structure interaction. A stabilised finite-element tsunami model
using tetrahedral meshes, volume-of-fluid interface capture and
SUPG/PSPG stabilisation has demonstrated
two-dimensional--three-dimensional tsunami coupling
\parencite{takase2016hybrid}. However, adopting this route would require
a separate finite-element fluid framework and additional treatment to
obtain the local face-conservative structure required by the present
architecture.

Structured staggered-grid projection methods provide strong
pressure--velocity coupling and efficient regular-grid operators. Their
use for the present barrier problem would require immersed-boundary or
cut-cell techniques for irregular geometry, together with a separate
Cartesian solver infrastructure.

Particle methods provide natural treatment of severe free-surface
deformation. However, pressure reconstruction, wall-traction transfer,
explicit representation of the air phase and systematic URANS--LES
operation remain less mature for the present requirements. The modified
moving-particle semi-implicit method reviewed in this study required
changes to the kernel, pressure-Poisson source and free-surface
detection to reduce pressure fluctuation sufficiently for
tsunami--structure comparison \parencite{huang2015modifiedmps}.

The alternative methods will therefore remain in the research register
as validation references and possible specialist methods, rather than as
the production local-fluid solver.

\paragraph{Decision Status and Remaining Numerical Choices.}
\label{sec:num-local-three-dimensional-decision-status}

\begin{table}[htbp]
  \centering
  \small
  \renewcommand{\arraystretch}{1.2}
  \caption{Decision status for the local three-dimensional fluid solver}
  \label{tab:num-local-three-dimensional-decision-status}
  \begin{tabular}{
      p{0.17\linewidth}
      p{0.28\linewidth}
      p{0.37\linewidth}
    }
    \hline
    Status
    &
    Component
    &
    Current decision
    \\
    \hline
    Selected
    &
    Spatial discretisation
    &
    Eulerian, cell-centred polyhedral finite-volume method
    \\
    Selected
    &
    Free-surface family
    &
    Volume-of-fluid interface capture
    \\
    Selected
    &
    Incompressibility treatment
    &
    Pressure-correction or projection family
    \\
    Selected
    &
    Turbulence architecture
    &
    Common solver supporting URANS $k$--$\omega$ SST for operational
    assessment and LES for selected comparison cases
    \\
    Selected
    &
    Software direction
    &
    OpenFOAM as the principal open-source implementation and
    architectural reference
    \\
    Provisional
    &
    Interface advection
    &
    Geometric volume-of-fluid method, with isoAdvector as the leading
    candidate
    \\
    Provisional
    &
    Interface pressure treatment
    &
    Sharp Ghost Fluid treatment rather than a purely diffuse mixture
    treatment
    \\
    Unresolved
    &
    Pressure correction
    &
    Exact PISO, PIMPLE or related correction sequence
    \\
    Unresolved
    &
    Moving structural boundary
    &
    Fixed terrain and fixed barrier patches are selected for the
    hydrodynamic baseline; mesh motion, immersed boundary, cut cell or
    overset treatment remains open only for activated FSI
    \\
    Unresolved
    &
    LES closure and wall treatment
    &
    LES sub-grid model, wall-function constants and rough-wall
    treatment
    \\
    Unresolved
    &
    Discretisation details
    &
    Temporal scheme, momentum convection, interface compression, linear
    solvers and preconditioners
    \\
    Unresolved
    &
    Mesh resolution
    &
    Static refinement, adaptive refinement and interface-focused mesh
    criteria
    \\
    \hline
  \end{tabular}
\end{table}

The next breadth-first task is to select the numerical family for the
structural equations of motion. Detailed choices within the local
three-dimensional fluid solver will be revisited after the structural
solver and the principal coupling frameworks have been selected.

\subsection{Unresolved Decisions}
\label{sec:res-num-spec-unresolved-decisions}

% TODO: Record unresolved questions, required evidence and decision owners.

The unresolved scientific and numerical decisions are:

\begin{itemize}
  \item final regional dry-depth tolerance and velocity regularisation
    constants;
  \item final regional HLL/HLLC wave-speed safeguards;
  \item final regional mesh levels and domain-specific sponge placement;
  \item exact local interface-advection algorithm;
  \item exact local pressure-correction sequence;
  \item exact local LES closure, wall-function constants and rough-wall
    treatment;
  \item local structural-boundary representation;
  \item local mesh-resolution and refinement criteria;
  \item structural element and time-integration method;
  \item numerical FSI coupling algorithm;
  \item numerical two-dimensional--three-dimensional transfer operators.
\end{itemize}

\subsection{Requirements and Handoffs}
\label{sec:res-num-spec-requirements}

% TODO: State requirements passed to other Research Domains, Software,
% verification activities or later execution work.

The next research pass shall close the following non-regional-model
items:

\begin{enumerate}
  \item QGIS/geospatial domain and dataset lock;
  \item regional verification and validation test execution;
  \item structural dynamics solver;
  \item two-dimensional--three-dimensional coupling solver;
  \item FSI coupling solver;
  \item software implementation details for the locked regional and
    local hydrodynamic specifications.
\end{enumerate}

Software work may begin with modular data structures and scaffolds:
finite-volume mesh and field containers, manifest-driven I/O,
boundary-face state containers, regional 2D scaffold, replay-boundary
format and local 3D scaffold. This Research deliverable does not
authorise implementation of the local URANS--VOF solver before the
remaining numerical, validation and data-interface decisions are
accepted.

\subsection{Deliverable Completion Record}
\label{sec:res-num-spec-completion-record}

% TODO: Record whether the Deliverable is:
%
% - Not started
% - In progress
% - Draft complete
% - Reviewed
% - Accepted
%
% Acceptance requires:
%
% - the research questions to be answered;
% - claims to be supported by appropriate evidence;
% - assumptions and limitations to be explicit;
% - the project-specific conclusion to be stated;
% - downstream requirements to be traceable;
% - unresolved decisions to be closed or formally deferred.

Regional 2D NLSWE and local 3D URANS--VOF numerical-model status:
research specification complete; implementation and validation pending.
